(1) 二维波矢态密度每面积是 \(d^2k/(2\pi)^2=kdk/(2\pi)\)。每模平均能 \(\hbar ck/(e^{\beta\hbar ck}-1)\)，所以 \(U/A=g\hbar c\int_0^\infty k^2dk/[2\pi(e^{\beta\hbar ck}-1)]=\boxed{g\zeta(3)(k_BT)^3/(\pi\hbar^2c^2)}\)，利用 \(\int_0^\infty x^2/(e^x-1)dx=2\zeta(3)\)。

(2) 玻色模自由能 \(F=gAk_BT\int d^2k/(2\pi)^2\ln(1-e^{-\beta\hbar ck})\)。分部积分给 \(F=-U/2\)，所以二维声学模热压强 \(P=-(\partial F/\partial A)_T=U/(2A)\)，是二维线性色散粒子能流各向同性的结果。

(3) 令 \(u_0=g\zeta(3)k_B^3/(\pi\hbar^2c^2)\)，有 \(U=A u_0T^3\)、\(F=-Au_0T^3/2\)。于是 \(S=-\partial_TF=(3/2)Au_0T^2\)、\(C_A=\partial_TU=3Au_0T^2\)，随 \(T\to0\) 均趋零；要求理想连续声学色散延续到相关热波长。

(4) 绝热可逆 \(S=(3/2)Au_0T^2\) 守恒，所以 \(\boxed{T\sqrt A=\mathrm{const}}\)。当最短热激发波长与样品尺度相当时离散模不能以连续态密度处理；温度过高时声速、极化数和色散将偏离线性假设。
